\section{Introduction}
\label{sec:intro}

Predicting the response of a physical system to prescribed control actions is an important part of transient analysis and system design. In nuclear and other energy systems a change in control may affect several coupled processes over different time scales. Reactor power may respond rapidly while temperatures continue to evolve long after the control action has ended. Physics-based simulations can resolve these relationships and provide information that is difficult to obtain experimentally. However the computational cost of repeated simulations can limit their use in the evaluation of many possible operating trajectories. Machine learning provides an alternative through surrogate models that approximate the response of a reference simulator. Such models must reproduce the transient behavior of the system and provide useful estimates of prediction uncertainty.

The development of an accurate transient surrogate begins with the training data. For a dynamical system the response at a given time depends on the preceding evolution of the physical state. Two control trajectories may therefore reach the same actuator position while producing different responses. The rate of motion and the duration of each operating condition can have a substantial influence on the resulting trajectory. Sampling individual control values is not sufficient to represent these temporal effects. Instead the dataset must contain complete control trajectories and the corresponding multivariate response trajectories produced by the simulator. This provides the model with examples of how changes in control propagate through the coupled system.

An additional challenge arises when the surrogate is used to forecast beyond its observed history. During one-step training a recurrent model receives input sequences containing reference responses. In autoregressive prediction the model instead uses its previous predictions to construct subsequent input sequences. Small errors can accumulate and influence the remainder of a forecast. High accuracy during one-step training does not necessarily imply accurate prediction over a complete transient. The problem becomes more difficult for multivariate forecasting because errors may grow at different rates across the output variables.

Predictive uncertainty is consequently an important part of transient surrogate modeling. Deep ensembles provide a practical way to measure disagreement among independently trained models. Their spread can vary with operating conditions and along a forecast trajectory. However ensemble disagreement alone does not establish the probability that a reference response will fall within a particular interval. Conformal prediction offers a complementary approach that uses held-out observations to determine interval widths with a specified marginal coverage objective. The construction of the calibration score determines whether coverage is assessed at an individual horizon or over an entire forecast trajectory.

This work develops the Multivariate Adaptive Conformal Ensemble (MACE) framework for uncertainty quantification in autoregressive multivariate forecasting. MACE combines the mean and spread of an LSTM ensemble with conformal calibration. The analysis begins with temporally correlated control-trajectory sampling and examines the accuracy of recursive ensemble prediction. We then investigate horizon-based and trajectory-based conformal intervals under different calibration distributions. Finally the method is evaluated using operation-informed transients from a reduced heat-pipe microreactor model. The objective is to assess predictive accuracy and interval reliability within a consistent forecasting framework rather than to optimize or implement a control policy.

\subsection{Literature review}
\label{sec:intro:literature}

The performance of a data-driven physical surrogate depends strongly on the simulations used for training. For nonlinear systems the available data determine which operating conditions are represented directly and which behaviors must be inferred. Pestourie \emph{et al.}\ introduced an active learning approach for surrogate modeling of partial differential equations and demonstrated its application to metasurface design~\cite{pestourie2020surrogates}. Kapadia \emph{et al.}\ investigated active data selection for nonlinear parametric systems with reduced representations~\cite{kapadia2024surrogate}. These studies demonstrate the importance of selecting informative simulations. They also show that data generation should be considered part of surrogate development rather than a separate preprocessing task.

Time-dependent physical systems introduce additional considerations because the temporal structure of an input affects the resulting response. Control trajectories with similar amplitudes may produce different transients when their rates or durations differ. Zimmer \emph{et al.}\ considered temporally structured inputs in their study of safe active learning for time-series models using Gaussian processes~\cite{zimmer2018safe}. In contrast to uncertainty-driven acquisition the present work uses a specified stochastic sampling procedure to generate control trajectories before model training. This distinction is important because the resulting dataset is defined by the control sampling distribution rather than by predictions from an evolving surrogate.

Gaussian processes provide a convenient mathematical representation for correlated temporal variations. The Mat\'ern covariance family allows the smoothness and correlation of a stochastic process to be controlled through a small number of parameters~\cite{porcu2024matern}. In this study the covariance is applied to control velocity. Integrating the sampled velocities produces control trajectories with correlated motion and an explicitly defined initial position. All original stochastic trajectories were drawn using the same velocity semivariogram. The retained profiles were specified to satisfy an angular velocity magnitude below $1^\circ\,\mathrm{s}^{-1}$. This requirement concerns the accepted data rather than a mathematical bound imposed by the covariance function. The effects of initialization and trajectory rejection must also be considered when characterizing the resulting distribution. A separate operation-informed sampling procedure is used later to evaluate control trajectories with different temporal characteristics.

Long short-term memory (LSTM) networks provide an established approach to learning relationships within sequential data. Their recurrent gates allow information from preceding time steps to influence future predictions~\cite{hochreiter1997lstm}. Radaideh \emph{et al.}\ applied neural time-series forecasting to loss-of-coolant accident transients in nuclear power plants~\cite{radaideh2020loca}. Appiah \emph{et al.}\ investigated LSTM-based system identification for the control of sodium cold-trap purification temperatures~\cite{appiah2024coldtrap}. These applications demonstrate the usefulness of sequence models for nuclear system dynamics. Nevertheless the forecasting problem depends on how the trained model is used. A model evaluated using reference histories is not equivalent to a model required to predict an entire transient from its own preceding outputs.

Autoregressive forecasting allows a trained recurrent model to predict an extended response without access to future observations. Following initialization each prediction becomes part of the input used for the next step. Errors in the estimated state may therefore influence later predictions even when the prescribed future control is known. This effect is particularly relevant for coupled physical systems in which an error in one output can affect the prediction of several others. Evaluating the model at individual steps may conceal this behavior. Complete forecast trajectories are needed to characterize accumulated error and to distinguish performance during control motion from the subsequent thermal response.

Ensemble learning provides a practical means of studying the variability of such predictions. Lakshminarayanan \emph{et al.}\ demonstrated the use of independently trained deep neural networks for predictive uncertainty estimation~\cite{lakshminarayanan2017deep}. Differences in model initialization and training data can lead ensemble members to produce different predictions. In autoregressive forecasting each member may also follow a different predicted history. The resulting spread can change substantially over time and across output variables. This makes ensemble disagreement a useful candidate for adaptive interval scaling. However the spread is not inherently a calibrated measure of predictive uncertainty and cannot be assigned a probability interpretation without further assessment.

The reliability of ensemble uncertainty estimates can also change when the model encounters data outside its original evaluation distribution. Ovadia \emph{et al.}\ investigated predictive uncertainty under dataset shift and demonstrated the importance of evaluating uncertainty in unfamiliar conditions~\cite{ovadia2019shift}. Their experiments focused primarily on classification and do not directly establish coverage for regression trajectories. Nevertheless the broader issue applies to physical surrogates. A model may retain reasonable mean prediction accuracy while its uncertainty intervals no longer achieve their intended coverage. Forecast accuracy and interval reliability must therefore be evaluated separately.

Conformal prediction provides a distribution-free approach to constructing prediction sets around an existing model. In split conformal regression the predictor is trained before the calibration scores are calculated. A quantile of those scores determines the size of the prediction interval. Under exchangeability between calibration and evaluation samples the resulting construction provides a finite-sample marginal coverage guarantee~\cite{lei2018regression,angelopoulos2023conformal}. The guarantee applies to the defined statistical sample and coverage event. It does not imply that the same coverage is achieved for every operating condition or at every point within an individual trajectory.

The construction of nonconformity scores provides considerable flexibility in how conformal intervals are formed. Papadopoulos \emph{et al.}\ studied normalized nonconformity measures that account for differences in prediction difficulty~\cite{papadopoulos2008normalized}. Romano \emph{et al.}\ introduced conformalized quantile regression to combine quantile prediction with calibration~\cite{romano2019cqr}. Both approaches allow interval width to vary with the forecast while retaining a conformal calibration step. MACE follows this general principle by using ensemble spread as an adaptive normalization factor. A positive additive constant ensures that the scale remains nonzero when the ensemble members produce nearly identical predictions. Calibration then determines the multiplier required for the chosen coverage objective.

The combination of ensembles and conformal prediction has also been investigated for temporal forecasting. Xu and Xie developed ensemble bootstrap prediction intervals (EnbPI) for time-series prediction~\cite{xu2021enbpi}. Their method combines ensemble forecasts with residual-based interval construction and addresses dependence within a sequential observation process. The present work considers a different statistical unit. Independent control-trajectory simulations provide individual calibration samples while observations within a trajectory remain temporally dependent. By calculating a score for each complete trajectory the conformal construction can account for dependence across its forecast horizons without treating every time step as an independent calibration observation.

A related question concerns the event that a prediction interval is expected to cover. Stankevi\v{c}i\={u}t\.{e} \emph{et al.}\ investigated conformal forecasting for time series~\cite{stankeviciute2021forecasting}. Sun and Yu examined multivariate multi-step conformal prediction through a copula-based approach~\cite{sun2024copulacpts}. These studies illustrate the importance of dependence across forecast coordinates. For a multivariate response the coverage of a single output at one horizon differs from coverage of that output over the entire trajectory. Requiring all output variables to be covered introduces a further distinction. These objectives should be defined directly in the calibration procedure rather than inferred from a reported average coverage percentage.

In particular horizon-specific calibration for each output does not ensure that all outputs are covered simultaneously at the same horizon. Likewise trajectory calibration performed separately for each output does not ensure that every output is covered throughout the complete forecast. MACE therefore distinguishes MACE-Horizon and MACE-Trajectory from an additional joint trajectory calibration. Horizon complete multivariate coverage and Trajectory complete multivariate coverage are also evaluated as separate empirical events. This distinction is necessary when interpreting intervals for a system whose outputs are physically coupled.

Conformal coverage also depends on the relationship between the calibration and evaluation distributions. If the model remains unchanged then its training data need not follow the calibration distribution. The relevant exchangeability condition concerns the calibration and evaluation samples used to assess the fitted predictor. However quantiles obtained from one population generally do not retain their nominal coverage guarantee under a different evaluation distribution. Barber \emph{et al.}\ studied conformal prediction beyond exchangeability and showed how departures from the standard assumptions require additional treatment~\cite{barber2023exchangeability}. In the present study recalibration on independent samples from the new operating distribution provides a direct comparison without changing the trained forecasting models.

Heat-pipe microreactor operation offers a useful application for this analysis. Price \emph{et al.}\ investigated reactivity control optimization using neural surrogates and evolutionary algorithms~\cite{price2022microreactor}. Tunkle \emph{et al.}\ examined control-drum operation during transient and load-following conditions using deep reinforcement learning~\cite{tunkle2025microreactor}. These studies motivate the use of surrogate models for predicting responses to drum motion. However neither provides a coverage guarantee for the ensemble and simulator considered here. The operation-informed control trajectories in our analysis contain prescribed motion and holding periods that differ from the original stochastic excitation. They therefore provide a useful setting for investigating the transfer of conformal intervals between control distributions.

The literature establishes methods for generating simulation data and learning time-dependent responses. It also provides tools for ensemble uncertainty estimation and conformal forecasting. The present work combines these elements with a defined trajectory-sampling procedure and an autoregressive multivariate surrogate. Particular attention is given to the relationship between ensemble spread and forecast error. We also examine how the choice of conformal score changes the coverage objective and whether recalibration improves interval reliability when the evaluation distribution changes. These questions form the basis of the methodology developed in the following section.

\subsection{Objectives and contributions}
\label{sec:intro:contributions}

This study introduces MACE as a framework for Multivariate Adaptive Conformal Ensemble forecasting. The approach combines stochastic control-trajectory sampling with an ensemble of autoregressive LSTM models and conformal prediction. Temporally correlated velocities are generated from a common Mat\'ern semivariogram and integrated into prescribed control trajectories. The resulting simulations provide multivariate response trajectories for model training and evaluation. Ten LSTM members are trained using separate subsets of the original training data. Each member then predicts the complete response using its own preceding predictions and the known control trajectory.

MACE uses the ensemble mean as the central forecast and the member spread as a local uncertainty scale. The spread is supplemented by a common positive additive floor to prevent vanishing denominators in the normalized nonconformity scores. Conformal calibration determines the interval multiplier using an independent set of trajectory samples. This construction allows interval width to adapt along the forecast while separating network training from uncertainty calibration. The method is compared with absolute-error conformal prediction under the same coverage objectives.

We consider two principal calibration methods. MACE-Horizon calculates a separate threshold for each output at each forecast horizon. MACE-Trajectory instead uses the largest normalized error along each output trajectory to calibrate a single threshold for that output. A third construction uses the maximum score across all outputs and horizons to target complete multivariate trajectory coverage. These methods differ in the event they seek to cover and consequently in their expected interval widths. We evaluate marginal coverage and complete multivariate coverage alongside full interval width and interval score.

The effect of calibration distribution is investigated using two populations of control trajectories. The first comprises the stochastic trajectories used to develop the surrogate. The second consists of operation-informed maneuvers characterized by sustained motion and holding periods. We compare intervals calibrated using the original stochastic population with intervals calibrated using independent operation-informed trajectories. The ensemble models and training scaler remain unchanged across the comparison. Recalibration therefore changes the prediction intervals without modifying their central forecasts.

The final application uses a reduced Modelica representation of a heat-pipe microreactor simulated in Dymola. The system couples point kinetics with lumped heat transfer through the reactor and balance of plant. Its secondary-side model includes a dynamic steam-generator metal node and calculated steam outlet conditions. The study evaluates offline forecasts against this reference simulation over 200~s. Results are considered both across complete trajectories and during the different periods of prescribed drum motion and thermal settling.

The primary contributions are summarized below.
\begin{enumerate}
 \item A temporally correlated control-trajectory sampling method based on a common velocity semivariogram and integrated velocity histories for generating admissible simulation inputs.
 \item An autoregressive ten-member LSTM forecasting approach that preserves individual member histories and correctly maps predicted outputs between the stored input and output scaling conventions.
 \item The Multivariate Adaptive Conformal Ensemble formulation with horizon-specific and trajectory-specific calibration together with a joint complete multivariate trajectory construction.
 \item An evaluation of forecasting accuracy and conformal calibration transfer using stochastic and operation-informed control distributions with matched forecasts and whole-trajectory statistical comparisons.
\end{enumerate}

The following section presents the control sampling procedure and the autoregressive ensemble formulation. It then describes MACE calibration and the reduced physical model before defining the evaluation procedure. The results examine prediction accuracy and conformal interval performance under the original stochastic distribution and after the change to operation-informed transients.
